\chapter{2016 Nobel Prize in Economics}
\textbf{Laureates: }Oliver Hart (USA/UK), Bengt Holmström (Finland)

\section*{Reason for Award}
Contract theory

\section{What They Did}
Oliver Hart and Bengt Holmström advanced contract theory, providing insights into how contractual arrangements should be designed to align incentives when parties have different information and objectives. Hart developed the incomplete contracts approach, showing that because contracts cannot specify all future contingencies, property rights and control rights matter. He analyzed how ownership structure affects investment incentives and firm boundaries. Holmström developed the principal-agent framework, showing how optimal contracts balance incentives with insurance when effort is unobservable. He analyzed how performance measures, pay-for-performance schemes, and career concerns affect agent behavior. Together, they created a comprehensive theory of contractual relationships in firms, markets, and organizations.

\section{Impact on Economic Thinking}
Hart and Holmström transformed understanding of contractual relationships and organizational design. Their work explained why firms have boundaries, how ownership affects investment incentives, and why contracts are often incomplete. This influenced corporate governance, executive compensation, and merger analysis. The incomplete contracts framework provided tools for analyzing public-private partnerships, outsourcing decisions, and joint ventures. Holmström's incentive theory informed compensation design in corporations, nonprofits, and public sector organizations. Hart's work on property rights clarified debates about privatization, regulation, and institutional design. Their combined contributions made contract theory a central pillar of modern economics, with applications spanning industrial organization, labor economics, public economics, and corporate finance.

\section{Modern Perspectives Today}
Contract theory remains essential for understanding modern economic organization. Hart's incomplete contracts framework informs debates about corporate governance, private equity, and firm boundaries. Holmström's incentive theory guides compensation design in the age of stock options, performance bonuses, and executive pay. Research on platform economies applies contract theory to understand relationships between platforms, providers, and consumers. The theory also informs labor market analysis, including gig economy arrangements and remote work contracts. Public policy debates about healthcare provision, education outsourcing, and infrastructure development draw on contract theory insights. As work arrangements evolve with technology and globalization, Hart and Holmström's theoretical tools remain indispensable for analyzing incentive compatibility and organizational design.
